\section{Methods}

\paragraph{Study design}
\label{sec:selecting}
We conducted a systematic review, as illustrated in  \Cref{fig:systematic-review-process}. Our corpus consisted of 46,114 articles drawn from the proceedings of ICML, ICLR and NeurIPS (accessed via proceedings websites)  between 2018 and 2024, and from ACL, NAACL and EMNLP between 2020 and 2024 (accessed via ACL Anthology). The ACL range was limited by abstract availability.

We identified and selected articles whose titles or abstracts contained the keywords `benchmark' and either `LLM' or `language model', resulting in an initial set of 2,189 articles, with most articles coming from recent years, and only 14 in 2018 and 2019.

We applied four inclusion criteria to assess the relevance and suitability of each article. First, we evaluated whether the article concerned the capabilities of LLMs, excluding those focused solely on technical aspects such as inference speed or energy consumption. Second, we determined whether the article introduced an empirical benchmark and reported LLM performance, excluding opinions, reviews or policy frameworks. Third, we assessed whether the benchmark was compatible with text and vision models, filtering out those that required other modalities such as audio or video. Finally, we checked that the article introduced a novel benchmark or made a substantial modification to an existing one, excluding repackaged or minimally altered combinations of prior benchmarks. 

%This criterion excludes articles that focus only on technical performance, such as computational efficiency or inference speed.

%     \item \textit{Empirical}: \textit{Does the article implement a benchmark and report the performance of large language models?} This criterion excludes opinions, reviews or policy frameworks.

%     \item \textit{Modality}: \textit{Could a text and vision model be tested on the benchmark?} This criterion excludes articles requiring video, audio, or other multimodal abilities that are not the focus of our review. 

%     \item \textit{Novelty}: \textit{Does the article introduce a new benchmark or meaningfully alter an existing one?} This criterion excludes articles which combine or repackage existing benchmarks with only minimal changes.
%\end{enumerate}

We first used GPT-4o mini \cite{gpt-4o-mini} to screen the articles on the basis of the first three criteria.  This model-assisted step was validated against human-labelled data for  a sample of 50 articles and achieved an F1 score of 84\%. This automated step reduced the set to 522 articles eligible for manual filtering. We then assigned the 522 eligible articles to 29 reviewers matched on area of expertise, to manually determine inclusion using all four criteria, resulting in 445 articles included for final review.


\paragraph{Codebook and expert review}\label{codebook}

We created an initial a priori codebook for phenomena, tasks, metrics and claims. Building on the definition of a benchmark from \textcite{rajiAIEverythingWhole2021a}, we consider a benchmark to be a `task' and `metric' which are used together to represent a `phenomenon' of interest. These elements are considered alongside the interpretation of the results by the authors.\footnote{We chose to use the terms `task' and `phenomenon', rather than `dataset' and `task', to better capture benchmarks that involve dynamic elements or aim to measure abstract capabilities rather than specific tasks.} For example, in GSM8K \cite{cobbe_training_2021}, the \textit{phenomenon} is `multi-step mathematical reasoning', which is measured via the \textit{task} of answering short free response questions drawn from grade-school mathematics word problems, which are scored via the `exact match' \textit{metric}.

Items in the codebook were derived deductively based on prior literature to provide indications of key aspects of construct validity, including face, predictive, content, ecological, convergent and discriminant validity (see~\Cref{background}). Each article was coded by a primary reviewer using this codebook. A second reviewer mapped the responses onto a simplified list of options for computing statistics and these mappings were verified by the primary reviewer. A random sample of 46 papers were reviewed twice, with a mean Brennan–Prediger Kappa of .524 across all 30 categorical questions.
%
%\paragraph{Analysis, Recommendations and Checklist} 
The first author then read a subset of 50 articles and reviewed all the 445 annotations to synthesise the findings into an initial set of recommendations through an inductive open coding process. Finally, these recommendations were collaboratively refined through an iterative process involving multiple authors across five meetings.
